% Part: first-order-logic
% Chapter: introduction
% Section: first-order-logic

\documentclass[../../../include/open-logic-section]{subfiles}

\begin{document}

\olfileid{fol}{int}{fol}

\olsection{一阶逻辑}

你可能在初次接触形式逻辑时就已经熟悉了一阶逻辑。\footnote{事实上，我们或多或少假设你熟悉它！如果你不熟悉，可以复习更基础的教科书，例如 \emph{forall x} \citep{Magnus2021}。}你可能知道它叫“量化逻辑”或“谓词逻辑”。一阶逻辑首先是一种形式语言。这意味着它有一套词汇表，其表达式是由该词汇表中的符号组成的字符串。但并非任意字符串都是合法的。合法的表达式有不同的种类：项、公式和语句。我们主要对一阶逻辑的语句感兴趣：它们为我们提供了英语句子的形式对应物，对于它们，我们可以提出逻辑学家通常感兴趣的问题。例如：

\begin{itemize}
    \item $!B$ 是否从 $!A$ 逻辑地得出？
    \item $!A$ 是逻辑真、逻辑假还是偶然的？
    \item $!A$ 和 $!B$ 是否等价？
\end{itemize}

这些问题主要是关于一阶逻辑语句“意义”的问题。例如，哲学家会这样分析“$!B$ 是否从 $!A$ 逻辑地得出”这个问题：是否存在 $!A$ 为真而 $!B$ 为假的情况（此时 $!B$ 不从 $!A$ 得出），还是说任何使 $!A$ 为真的情况也都使 $!B$ 为真（此时 $!B$ 确实从 $!A$ 得出）？但我们还不清楚“情况”究竟指什么——这正是\emph{语义学}的任务。一阶逻辑的语义学为哲学家关于“情况”的直观概念提供了一个数学上精确的模型，并且同样重要的是，也为语句~$!A$ 在一个情况中\emph{为真}意味着什么提供了精确的模型。我们将要构造的这个数学上精确的模型称为结构。使“在其中为真”这一概念精确化的关系，被称为\emph{满足}关系。因此，我们将要定义的是：对于语句~$!A$ 和结构~$\Struct{M}$，“$!A$ 在结构~$\Struct{M}$ 中被满足”（记作 $\Sat{M}{!A}$）。一旦完成这个定义，我们也可以给出其他语义术语（如“从……得出”或“逻辑真”）的精确定义。有了这些定义，我们就能以数学的精确性来判定，例如，$\lforall[x][(!A(x) \lif !B(x)),
\lexists[x][!A(x)] \Entails \lexists[x][!B(x)]]$是否成立。答案当然是“是”。如果你接受过用一阶逻辑将英语句子符号化的训练，你会认出这正是诸如“所有蚂蚁都是昆虫，有蚂蚁，因此有昆虫”这类句子的符号化。这显然是一个有效的论证，因此我们为形式语言建立的“从……得出”的数学模型也应该给出相同的答案。

你可能还记得形式逻辑入门中的另一个主题：\emph{推导}。如果你上过一门形式逻辑入门课，你的老师很可能让你练习过寻找这样的推导，也许正是寻找一个表明上述语义后承关系成立的推导。给出推导有很多不同的方法：你可能学过被称为“自然演绎”或“真值树”的方法，但还有很多其他方法。推导系统的目的是提供工具，用以回答上面提到的逻辑学家的问题：例如，一个以$\lforall[x][(!A(x) \lif !B(x))$和$\lexists[x][!A(x)]$为前提、以$\lexists[x][!B(x)]]$为结论（最后一行）的自然演绎推导，就\emph{验证}了$\lexists[x][!B(x)]$从$\lforall[x][(!A(x) \lif !B(x))]$和$\lexists[x][!A(x)]$逻辑地得出。

但这是为什么呢？表面上看，推导系统与语义学毫无关系：给出一个形式推导仅仅涉及以某种遵循规则的方式排列符号；它们根本不涉及“情况”或“在其中为真”。推导系统与语义学之间的联系必须通过元逻辑研究来建立。我们需要的是一个数学证明，例如，证明能够构造出从前提$\lforall[x][(!A(x) \lif !B(x))]$和$\lexists[x][!A(x)]$到$\lexists[x][!B(x)]$的形式推导，当且仅当$\lforall[x][(!A(x)
\lif !B(x))$和$\lexists[x][!A(x)]$共同语义后承$\lexists[x][!B(x)]]$。然而，在此之前，必须进行大量艰苦的工作，以正确地给出语法和语义的定义。

\end{document}